\clearpage
\section*{ANEXOS}
\addcontentsline{toc}{section}{ANEXOS}

\subsection*{Premissas de leitura dos resultados}

A análise utiliza quatro naturezas de informação: dados publicados de fabricantes e instituições; evidências qualitativas de campo; cálculos sobre parâmetros declarados; e critérios de aceitação do produto. A Tabela~\ref{tab:premissas-analise} reúne as condições de interpretação dos principais resultados.

\begin{table}[H]
\caption{Premissas técnicas e econômicas da análise}
\label{tab:premissas-analise}\centering\footnotesize
\begin{tabularx}{\textwidth}{L{3.1cm}YL{4.4cm}}
\toprule\textbf{Tema}&\textbf{Condição considerada}&\textbf{Natureza e alcance}\\\midrule
Campo&Uma obra residencial; entrevistas com um operário e um distribuidor&Caracterização qualitativa do uso, sem taxa estatística de avaria.\\
Mix&Uma cuba ou corpo individual de bacia por dispositivo; kit montado e monobloco fora da faixa&Envelope de aplicação e limites de projeto apoiados em catálogos.\\
Reserva dimensional&$\pm10$\,mm sobre as faixas nominais&Margem de engenharia; a tolerância aproximada publicada pela Celite não define todas as louças.\\
Carga e espaço&40\,kg de louça; 8\,kg de dispositivo; base 680\,$\times$\,480\,mm&Orçamento de desenvolvimento; cálculo geométrico e estático.\\
Empilhamento&Dois níveis por 24\,h; compressão-meta de 7,5\,kN&Condição de avaliação e critério de aceitação, sem carga nominal homologada.\\
Custos de cenário&R\$\,15 de embalagem; R\$\,500 de peça; R\$\,120 de reentrega; taxas de 0,5\% a 5\%&Análise paramétrica; não é cotação ou frequência observada no distribuidor.\\
Seleção funcional&M2 datum; 23 critérios; pesos 0--5; notas $-1$/0/$+1$&Julgamento de fronteira e capacidades, distinto de ensaio ou homologação.\\
Modularização&C1 como configuração exploratória; escala 0/1/3/5; OP-01 fora da base&Referência de motivos de separação, sem seleção física ou certificação de interfaces.\\
\bottomrule\end{tabularx}
\fonte{Síntese das condições de aplicação, dos cálculos e dos critérios apresentados no relatório.}
\end{table}
